% Recovery of 02_trit_geometrias and of holonomy with memory.
\subsubsection{Temporal parametrization of compatibility and tritic orientation}
\label{trit-temporal-compatibilidad}

The temporal interpretation of TRIT is established on transitions of the
complete state. Let \(x\xrightarrow{\rm adm}y\) be the admissibility relation determined
by the TPK operations on \(\widehat X\). A path parametrized
by a discretely ordered set \(\mathcal T\) is a map
\[
 \gamma:\mathcal T\longrightarrow\widehat X,\qquad
 \gamma(t)\xrightarrow{\rm adm}\gamma(t^+),
\]
where \(t^+\) is the successor of \(t\). The compatibility relation precedes
the choice of parameter: the order makes it possible to express its transitions as
an evolution. When the term section is used, the projection
\(p:E\to\mathcal T\) and the condition
\(p\circ\gamma=\operatorname{id}_{\mathcal T}\) are also specified. This determines
the state space, the temporal order, and the information being transported.

The signature \(\tau\) selects a quadratic regime. Relative to that
parametrization, it also admits a temporal reading. Type \(+1\), with
\(J_{+1}^{2}=-I\), represents return and memory; type \(0\), with
\(J_0^{2}=0\), represents the transition threshold; type \(-1\), with
\(J_{-1}^{2}=I\), represents opening along two eigendirections. In the

temporal terminology of the model, these functions correspond,
respectively, to past, present, and future. The correspondence concerns
the operational position of a transition within a history:
preserved antecedent, update, and admissible extension.

This interpretation acquires content through three distinct operations.
Restriction of a path recovers the prefix already constructed. The
update \(\mathcal U_t\) composes selection, transport, and recording at
the observed transition. Extension considers the continuations that
preserve the boundary conditions of that prefix. If
\(\mathcal P_n(x)\) denotes the set of admissible continuations of
length \(n\) from \(x\), deletion of the final step defines
\[
 r_n:\mathcal P_{n+1}(x)\longrightarrow\mathcal P_n(x).
\]
An unbounded continuation is a collection \((\gamma_n)_{n\geq1}\) satisfying
\(r_n(\gamma_{n+1})=\gamma_n\). The inverse system thus preserves
the constraints that each depth imposes on the preceding ones.
The existence of a compatible collection is established in the corresponding
extension domain; merely defining the sets
\(\mathcal P_n(x)\) does not replace that proof.

Bidirectionality therefore combines reconstruction of antecedents
with compatibility of continuations. On the sequence of signatures, the
involution \(C_{\rm rev}\) reverses the order and changes the sign of every
component. The position, sheet, and boundary transformations specified by
the transport also act on the path. The reversed signature is an observation
of the conjugate path, not a replacement for its remaining coordinates.
The threshold remains distinguished: the visible value \(0\) retains its
transition function even when its quotient or historical record is nonzero.

The relation to exponential geometries is equally precise. In
each fixed regime, \(\exp(tJ_\tau)\) composes parameters and preserves the
algebraic norm. Elliptic closure belongs to that local representation.
The phase return of \(\Gamma_9\), by contrast, acts on the path
with memory, in accordance with \eqref{eq:holonomia-memoria}. A path can
recover its phase while preserving a new antecedent: the historical datum has
changed although the phase observation is the same. The parabolic threshold
retains a linear displacement; the hyperbolic opening separates two
eigendirections of expansion and contraction. These three realizations
provide composition laws, while the TPK determines their effective
selection and concatenation.

The aperiodic suspension of Section \ref{sec:generacion} will make
this distinction explicit through a finite phase factor, an irrational rotation,
and a memory degree. The link between TRIT and temporality is thus
expressed by operations on paths and by their quadratic
representations. Its content goes beyond a classification of signs: it preserves the
regime, the direction of transport, and the relation among antecedent,
transition, and continuation.
